% Pursuing Stacks (cotes 134-2 to 134-8): archive cote 134-3, site page 80 (PDF page 81), typescript page 80.
% typescript. Draft compared: ps2.tex:4144-4223.
% AI-transcribed by Claude Opus 5.5 (claude-opus-5-5, adaptive thinking, effort high), Message Batches API, batch ps1 (msgbatch_01Rv5hCaNTtUt4JHx1yqoHPL), 2026-11-12; not yet checked by eye.
\scanpage{134-3/80}
\mspage{(80)}

\margin{\textcircled{}}(iii) As in (ii), but moreover the adjunction morphism in $W_A^{-1}A\hat{}$ for
the pair of quasi-inverse equivalences\inserted{(5)}, being \hand{likewise deduced from (4).}
(iv) For any C in (Cat), (3) is a weak equivalence.
\margin{\textcircled{}}(v) Same as (iv), with C restricted to having a final element, i.e.
for any such C, $j_Ai_A(C) = A_{/j_A(C)}$ is aspheric.
(vi) For any F in $A\hat{}$, (4) is a weak equivalence, moreover
\[ W_{(Cat)} = j_A^{-1}(W_A) \qquad , \]
i.e. a \xout{functor} map $f\colon C' \to C$ in (Cat) is a \xout{quasi}\inserted[typed]{weak}equivalence,
iff $j_A(f)$ is a \xout{quasi}\inserted[typed]{weak}equivalence in $A\hat{}$ - which means, by defi-%
nition (more or less) that $i_Aj_A(f)$ is a \xout{quasi}\inserted[typed]{weak}equivalence.
(vii) The functor
(6) \[ i_Aj_A \; : \; (Cat) \longrightarrow (Cat) \quad , \quad C \mapsto A_{/j_A(C)} \, , \]
transforms \xout{quasi}\inserted[typed]{weak}equivalences into \xout{quasi}\inserted[typed]{weak}equivalences , i.e.
gives rise to a functor (Hot)$\to$(Hot), and moreover \deleted{A is} \hand{the}
\deleted{aspheric.} \hand{latter respects final object (i.e. A is aspheric).}
(viii) Same as (vii), but restricting to \xout{quax} weak equivalences of the
type $C \to e$, where C has a final element and e is the final
object in (Cat) (the one-point category)\hand{, plus asphericity of A.}

The trivial implications between all these conditions can be summari\xout{z}
zed in the diagram
\diagram{\begin{tikzcd}[row sep=small, arrows=Rightarrow]
(iii) \ar[r,Leftrightarrow] \ar[d] & (vi) \ar[r] & (vii) \ar[r] & (viii) \ar[dl] \\
(ii) \ar[r] \ar[d] & (iv) \ar[r] & (v) & \\
(i) & , & &
\end{tikzcd}}
the only \inserted[typed]{slightly} less obvious implication here is (viii)\hand{$\Rightarrow$}(v), which is seen
by looking at the commutative square deduced from $C \to e$ (C in (Cat) as in
(v)\xout{\unsure{,}} namely with final element) by applying (3)
\diagram{\begin{tikzcd}
i_Aj_A(C) \ar[r] \ar[d] & C \ar[d] \\
A{=}i_Aj_A(e) \ar[r] & e
\end{tikzcd} \quad ,}
by assumption the vertical arrows are weak equivalences, and so is $A \to e$
(because A is supposed to be aspheric), therefore the same holds for the
fourth arrow left. On the other hand an easy asphericity argument
showed us \inserted{[3.30]} that (v) \hand{$\Longrightarrow$}(iii), hence all conditions (ii) to (viii) are
equivalent, and they are equivalent to the stronger form of (iv), \xout{that}
say (iv'), saying that (3) is \uline{aspheric} for any C in (Cat). The only
equivalence which is not quite clear yet is that (i) implies the other
conditions\hand{.} \deleted{under the formally weakest form (v) let's say.} But it is so
